\textbf{Phần 1. Trắc nghiệm lựa chọn}
\setcounterpageref{ex}{0}

\begin{ex} %[3D4N28]
Một bài toán giải bằng hai bước tính thường yêu cầu
\choice
{Thực hiện duy nhất một phép tính}
{Chỉ cần vẽ hình}
{Không cần thực hiện phép tính}
{\True Tìm một giá trị trung gian trước khi tính kết quả cuối cùng}
\loigiai{
Bài toán giải bằng 2 bước tính cần tìm bước trung gian trước khi tính đáp số.\\
\textbf{Đáp án của câu hỏi là D}
}
\end{ex}

\begin{ex} %[3D4N28]
Thùng thứ nhất có $18\text{ l}$ dầu, thùng thứ hai có nhiều hơn thùng thứ nhất $6\text{ l}$ dầu. Bước đầu tiên cần tìm là
\choice
{Tổng số dầu hai thùng}
{\True Số dầu ở thùng thứ hai}
{Hiệu số dầu}
{Số dầu thùng thứ nhất}
\loigiai{
Bước 1: Tính số dầu thùng thứ hai ($18 + 6 = 24\text{ l}$).\\
\textbf{Đáp án của câu hỏi là B}
}
\end{ex}

\begin{ex} %[3D4H28]
Hàng thứ nhất có 8 cây, hàng thứ hai có số cây gấp đôi hàng thứ nhất. Cả hai hàng có số cây là
\choice
{$16$ cây}
{$28$ cây}
{$20$ cây}
{\True $24$ cây}
\loigiai{
Hàng 2: $8 \times 2 = 16$ cây.\\
Cả hai hàng: $8 + 16 = 24$ cây.\\
\textbf{Đáp án của câu hỏi là D}
}
\end{ex}

\begin{ex} %[3D4H28]
Bạn Mai có 15 bông hoa, bạn Lan có ít hơn bạn Mai 3 bông hoa. Cả hai bạn có số bông hoa là
\choice
{\True $27$ bông}
{$25$ bông}
{$30$ bông}
{$24$ bông}
\loigiai{
Lan có: $15 - 3 = 12$ bông.\\
Cả hai bạn: $15 + 12 = 27$ bông.\\
\textbf{Đáp án của câu hỏi là A}
}
\end{ex}

\begin{ex} %[3D4H28]
Một cửa hàng buổi sáng bán được $24\text{ kg}$ gạo, buổi chiều bán giảm đi 3 lần so với buổi sáng. Cả ngày cửa hàng bán được số ki-lô-gam gạo là
\choice
{$30\text{ kg}$}
{\True $32\text{ kg}$}
{$34\text{ kg}$}
{$36\text{ kg}$}
\loigiai{
Chiều: $24 : 3 = 8\text{ kg}$.\\
Cả ngày: $24 + 8 = 32\text{ kg}$.\\
\textbf{Đáp án của câu hỏi là B}
}
\end{ex}

\begin{ex} %[3D4H28]
Có 3 bao gạo, mỗi bao nặng $15\text{ kg}$. Người ta đã dùng hết $20\text{ kg}$. Số gạo còn lại là
\choice
{$22\text{ kg}$}
{$30\text{ kg}$}
{\True $25\text{ kg}$}
{$35\text{ kg}$}
\loigiai{
Tổng: $15 \times 3 = 45\text{ kg}$.\\
Còn lại: $45 - 20 = 25\text{ kg}$.\\
\textbf{Đáp án của câu hỏi là C}
}
\end{ex}

\begin{ex} %[3D4V28]
Một lớp học có 16 học sinh nam, số học sinh nữ nhiều hơn số học sinh nam 4 bạn. Lớp học đó có tất cả bao nhiêu học sinh?
\choice
{\True $36$ học sinh}
{$32$ học sinh}
{$38$ học sinh}
{$40$ học sinh}
\loigiai{
Nữ: $16 + 4 = 20$ bạn.\\
Cả lớp: $16 + 20 = 36$ học sinh.\\
\textbf{Đáp án của câu hỏi là A}
}
\end{ex}

\begin{ex} %[3D4V28]
Mẹ mua 4 hộp bánh, mỗi hộp có 6 cái bánh. Mẹ chia đều số bánh đó cho 3 anh em. Mỗi người nhận được số cái bánh là
\choice
{$7$ cái}
{$9$ cái}
{\True $8$ cái}
{$6$ cái}
\loigiai{
Tổng số bánh: $4 \times 6 = 24$ cái.\\
Mỗi người nhận: $24 : 3 = 8$ cái bánh.\\
\textbf{Đáp án của câu hỏi là C}
}
\end{ex}

\begin{ex} %[3D4V28]
Đoạn thẳng thứ nhất dài $18\text{ cm}$, đoạn thẳng thứ hai dài gấp 3 lần đoạn thẳng thứ nhất. Đoạn thẳng thứ hai dài hơn đoạn thẳng thứ nhất bao nhiêu xăng-ti-mét?
\choice
{\True $36\text{ cm}$}
{$32\text{ cm}$}
{$40\text{ cm}$}
{$54\text{ cm}$}
\loigiai{
Đoạn 2: $18 \times 3 = 54\text{ cm}$.\\
Đoạn 2 hơn đoạn 1: $54 - 18 = 36\text{ cm}$.\\
\textbf{Đáp án của câu hỏi là A}
}
\end{ex}

\begin{ex} %[3D4V28]
Một trang trại có 48 con gà, số vịt bằng $\dfrac{1}{4}$ số gà. Hỏi trang trại nuôi tất cả bao nhiêu con gà và vịt?
\choice
{$56$ con}
{$64$ con}
{\True $60$ con}
{$58$ con}
\loigiai{
Số vịt: $48 : 4 = 12$ con.\\
Tổng: $48 + 12 = 60$ con.\\
\textbf{Đáp án của câu hỏi là C}
}
\end{ex}

\begin{ex} %[3D4V28]
Có 5 túi bi, mỗi túi có 12 viên bi. Người ta lấy ra 18 viên bi để chơi. Số viên bi còn lại là
\choice
{$40$ viên}
{$45$ viên}
{$48$ viên}
{\True $42$ viên}
\loigiai{
Tổng: $12 \times 5 = 60$ viên.\\
Còn lại: $60 - 18 = 42$ viên bi.\\
\textbf{Đáp án của câu hỏi là D}
}
\end{ex}

\begin{ex} %[3D4C28]
Hai thùng có tất cả $80\text{ l}$ dầu. Thùng thứ nhất có $32\text{ l}$ dầu. Hỏi thùng thứ hai có nhiều hơn thùng thứ nhất bao nhiêu lít dầu?
\choice
{$14\text{ l}$}
{\True $16\text{ l}$}
{$18\text{ l}$}
{$20\text{ l}$}
\loigiai{
Thùng 2: $80 - 32 = 48\text{ l}$.\\
Thùng 2 hơn thùng 1: $48 - 32 = 16\text{ l}$.\\
\textbf{Đáp án của câu hỏi là B}
}
\end{ex}

\textbf{Phần 2. Trắc nghiệm đúng sai}
\setcounterpageref{ex}{0}

\begin{ex} %[3D4H28]
Bác Ba nuôi 14 con thỏ trắng, số thỏ xám gấp đôi số thỏ trắng. Xét tính đúng sai:
\choiceTF
{\True Số thỏ xám bác Ba nuôi là 28 con.}
{\True Tổng số thỏ trắng và thỏ xám là 42 con.}
{Số thỏ xám nhiều hơn số thỏ trắng là 12 con.}
{\True Số thỏ xám nhiều hơn số thỏ trắng là 14 con.}
\loigiai{
\begin{itemize}
    \item Thỏ xám: $14 \times 2 = 28$ con. Đúng.
    \item Tổng: $14 + 28 = 42$ con. Đúng.
    \item Hiệu: $28 - 14 = 14$ con $\neq 12$. Ý 3 Sai, Ý 4 Đúng.
\end{itemize}
\begin{itemize}
    \item a) \textbf{Đúng}.
    \item b) \textbf{Đúng}.
    \item c) \textbf{Sai}.
    \item d) \textbf{Đúng}.
\end{itemize}
}
\end{ex}

\begin{ex} %[3D4V28]
Một cuộn dây dài $90\text{ m}$. Lần 1 cắt đi $\dfrac{1}{3}$ cuộn dây, lần 2 cắt đi $25\text{ m}$. Xét tính đúng sai:
\choiceTF
{\True Lần 1 cắt được $30\text{ m}$ dây.}
{\True Cả hai lần cắt được tất cả $55\text{ m}$ dây.}
{Cuộn dây còn lại $40\text{ m}$.}
{\True Cuộn dây còn lại $35\text{ m}$.}
\loigiai{
\begin{itemize}
    \item Lần 1: $90 : 3 = 30\text{ m}$. Đúng.
    \item Đã cắt: $30 + 25 = 55\text{ m}$. Đúng.
    \item Còn lại: $90 - 55 = 35\text{ m} \neq 40\text{ m}$. Ý 3 Sai, Ý 4 Đúng.
\end{itemize}
\begin{itemize}
    \item a) \textbf{Đúng}.
    \item b) \textbf{Đúng}.
    \item c) \textbf{Sai}.
    \item d) \textbf{Đúng}.
\end{itemize}
}
\end{ex}

\textbf{Phần 3. Trả lời ngắn}
\setcounterpageref{ex}{0}

\begin{ex} %[3D4H28]
Thùng 1 có $25\text{ l}$ nước, thùng 2 có nhiều hơn thùng 1 là $10\text{ l}$. Hỏi cả hai thùng có bao nhiêu lít nước?
\shortans{60}
\loigiai{
\textbf{Đáp số:} $60$\\
Thùng 2: $35\text{ l}$. Tổng: $25 + 35 = 60\text{ l}$.
}
\end{ex}

\begin{ex} %[3D4H28]
Bạn An có 12 cái kẹo, bạn Bình có số kẹo gấp 3 lần bạn An. Hỏi bạn Bình có nhiều hơn bạn An bao nhiêu cái kẹo?
\shortans{24}
\loigiai{
\textbf{Đáp số:} $24$\\
Bình: $36$ cái. Hiệu: $36 - 12 = 24$ cái.
}
\end{ex}

\begin{ex} %[3D4V28]
Có 4 bao gạo, mỗi bao nặng $18\text{ kg}$. Người ta đã bán được $40\text{ kg}$. Tính số gạo còn lại theo ki-lô-gam.
\shortans{32}
\loigiai{
\textbf{Đáp số:} $32$\\
Tổng: $72\text{ kg}$. Còn lại: $72 - 40 = 32\text{ kg}$.
}
\end{ex}

\begin{ex} %[3D4V28]
Một lớp có 32 học sinh, trong đó $\dfrac{1}{4}$ là học sinh giỏi, còn lại là học sinh khá. Hỏi lớp đó có bao nhiêu học sinh khá?
\shortans{24}
\loigiai{
\textbf{Đáp số:} $24$\\
Giỏi: $32 : 4 = 8$ bạn. Khá: $32 - 8 = 24$ bạn.
}
\end{ex}

\begin{ex} %[3D4V28]
Mẹ mua 3 vỉ trứng, mỗi vỉ có 10 quả trứng. Mẹ đã dùng hết 14 quả. Hỏi mẹ còn lại bao nhiêu quả trứng?
\shortans{16}
\loigiai{
\textbf{Đáp số:} $16$\\
Tổng: $30$ quả. Còn lại: $30 - 14 = 16$ quả.
}
\end{ex}

\begin{ex} %[3D4C28]
Một cửa hàng có $75\text{ m}$ vải. Buổi sáng bán được $15\text{ m}$, buổi chiều bán được số vải gấp đôi buổi sáng. Hỏi cửa hàng còn lại bao nhiêu mét vải?
\shortans{30}
\loigiai{
\textbf{Đáp số:} $30$\\
Chiều: $30\text{ m}$. Đã bán: $15 + 30 = 45\text{ m}$. Còn lại: $75 - 45 = 30\text{ m}$.
}
\end{ex}
